\newpage
\chapter{PENDAHULUAN} \label{Bab I}

\section{Latar Belakang} \label{I.Latar Belakang}
Diagnosis penyakit kulit berbasis citra medis kini bertransisi secara signifikan ke arah paradigma "Deployment-First Medical AI", di mana efisiensi komputasi dan kecepatan inferensi pada perangkat edge menjadi prioritas utama selain akurasi klinis \cite{nguyen2026efficient, mehta2022mobilevit}. Di fasilitas pelayanan kesehatan dengan infrastruktur komputasi terbatas, ketersediaan model kecerdasan buatan yang ringan dan andal sangat krusial untuk membantu klinisi dalam melakukan deteksi dini lesi kulit \cite{lavaei2025resource}.\par

Dalam implementasinya, klasifikasi penyakit kulit berbasis Deep Learning dihadapkan pada dua tantangan utama: ketidakseimbangan kelas (class imbalance) yang ekstrem dan keterbatasan sumber daya komputasi \cite{esteva2017dermatologist, litjens2017survey}. Penelitian terdahulu pada dataset DermaMNIST menggunakan model konvensional seperti ResNet-18 dan ResNet-50 \cite{he2016deep} menunjukkan kecenderungan prediksi yang sangat terbias pada kelas mayoritas (Melanocytic Nevi mendominasi hingga $\sim$67\% data), sehingga gagal mengenali kelas-kelas penyakit langka secara akurat \cite{hariz2026analisis}. Di sisi lain, meskipun arsitektur Vision Transformer (ViT) \cite{dosovitskiy2021image, shamshad2022transformers} menawarkan keunggulan dalam menangkap konteks spasial global, kompleksitas komputasi kuadratiknya $\mathcal{O}(N^2)$ terhadap panjang sekuens/resolusi gambar sangat membebani memori VRAM dan daya komputasi perangkat lunak medis \cite{amangeldi2025cnn}.\par

Sebagai terobosan mutakhir, muncul arsitektur berbasis State Space Models (SSM) / Mamba yang menggabungkan lapisan konvolusi untuk fitur lokal dengan selective scan module untuk pemrosesan konteks global ber-kompleksitas linear $\mathcal{O}(N)$ \cite{gu2024mamba}. Salah satu arsitektur pionir di domain medis adalah MedMamba \cite{yue2024medmamba}, yang hadir dalam berbagai skala parameter meliputi MedMamba-T (Tiny, $\sim$14,5M parameter), MedMamba-S (Small, $\sim$22,8M parameter), dan MedMamba-B (Base, $\sim$47,1M parameter). Kehadiran variasi skala ini memungkinkan eksplorasi trade-off antara kapasitas representasi fitur dan efisiensi komputasi \cite{aalishah2025medmambalite}.\par

Namun demikian, validitas ilmiah dari pengembangan model AI medis sangat bergantung pada integritas data yang digunakan. Investigasi kualitas data oleh Abhishek et al. (2025) mengungkapkan ditemukannya kebocoran data (\textit{data leakage}) sebesar \textbf{13,47\%} (1.006 dari 7.470 lesi unik) pada dataset DermaMNIST murni akibat pemisahan data yang tidak terisolasi di tingkat pasien/lesi (\textit{lesion-level split}) \cite{abhishek2025investigating, yang2023medmnist}. Kebocoran ini menyebabkan performa benchmark terinflasi secara semu karena model sekadar "menghafal" lesi pasien yang sama di folder pengujian. Masalah serupa berupa duplikasi dan ketidaksesuaian label juga ditemukan pada dataset Fitzpatrick17k \cite{groh2021evaluating}. Oleh karena itu, penggunaan versi dataset terkoreksi (\textbf{DermaMNIST-C} dan \textbf{Fitzpatrick17k-C}) yang bebas dari kebocoran data sangat mutlak diperlukan.\par

Untuk membangun model diagnostik yang tangguh (\textit{robust}) \cite{disalvo2024medmnistc}, pengujian tidak hanya perlu dilakukan secara \textit{in-dataset}, melainkan juga melalui \textbf{evaluasi lintas-dataset (\textit{cross-dataset generalization})}. Selain itu, paradigma pelatihan juga berkembang dari \textit{Supervised Learning} langsung menuju \textit{Contrastive Learning} untuk mempelajari representasi visual lesi yang lebih kaya tanpa terdistorsi oleh bias label mayoritas \cite{khosla2020supervised}. \par

Berdasarkan urgensi tersebut, penelitian ini dirancang untuk melakukan evaluasi komprehensif terhadap \textbf{skalabilitas dan generalisasi arsitektur MedMamba (MedMamba-T, MedMamba-S, dan MedMamba-B)} pada dataset terkoreksi (DermaMNIST-C dan Fitzpatrick17k-C). Penelitian ini membandingkan skenario pelatihan \textit{Supervised} vs \textit{Contrastive}, menguji dampak \textit{data leakage} secara eksperimental, serta mengeksekusi \textbf{studi ablasi mendalam pada MedMamba-S} sebagai model fokus utama.\par

\section{Rumusan Masalah} \label{I.Rumusan Masalah}
Berdasarkan latar belakang yang telah dipaparkan, rumusan masalah dalam penelitian ini adalah sebagai berikut: \par

\begin{enumerate}[noitemsep]
	\item Bagaimana efisiensi komputasi, performa klasifikasi, dan skalabilitas dari keluarga arsitektur MedMamba (MedMamba-T, MedMamba-S, dan MedMamba-B) dalam mengklasifikasikan citra penyakit kulit dengan ketidakseimbangan kelas yang ekstrem? 
	\item Sejauh mana perbedaan efektivitas skenario pelatihan Supervised Learning dibandingkan Contrastive Learning dalam mengekstrak representasi fitur visual lesi kulit yang diskriminatif?
	\item Bagaimana tingkat keandalan generalisasi lintas-dataset (cross-dataset generalization) dari model MedMamba ketika dilatih dan diuji secara silang antara DermaMNIST-C dan Fitzpatrick17k-C, serta seberapa besar dampak kebocoran data (data leakage) yang terjadi pada dataset DermaMNIST murni (origin)?
	\item Sejauh mana komponen-komponen utama pada arsitektur fokus MedMamba-S (skema fungsi kerugian/loss, strategi inisialisasi bobot pretrained, dan augmentasi data) berkontribusi terhadap peningkatan performa klasifikasi melalui studi ablasi?
\end{enumerate}

\section{Tujuan Penelitian} \label{I.Tujuan}

Berdasarkan rumusan masalah yang telah ditetapkan, tujuan penelitian ini adalah sebagai berikut: \par

\begin{enumerate}[noitemsep]
	\item Mengimplementasikan dan mengevaluasi tiga skala arsitektur MedMamba (MedMamba-T, MedMamba-S, dan MedMamba-B) guna memetakan trade-off antara jumlah parameter, penggunaan VRAM, dan akurasi diagnostik. 
	\item Menerapkan dan menganalisis perbandingan paradigma pelatihan berbasis Supervised Learning dan Contrastive Learning untuk penanganan ketidakseimbangan kelas medis.
	\item Menguji kemampuan generalisasi model melalui pengujian lintas-dataset (cross-dataset) serta mengkuantifikasi dampak kebocoran data dengan membandingkan evaluasi pada DermaMNIST-origin, DermaMNIST-C, dan Fitzpatrick17k-C.
	\item Melaksanakan studi ablasi terstruktur pada arsitektur MedMamba-S untuk mengisolasi efek dari skema loss function (Class-Weighted Loss), pre-training initialization, dan strategi augmentasi.
\end{enumerate}

\section{Batasan Masalah} \label{I.Batasan}
Berikut adalah batasan yang diberlakukan dalam penelitian ini: \par

\begin{enumerate}[noitemsep]
    \item Arsitektur model yang dievaluasi berfokus pada keluarga MedMamba (MedMamba-T, MedMamba-S, dan MedMamba-B).
    \item Dataset yang digunakan terbatas pada citra dermatologi publik: DermaMNIST-origin, DermaMNIST-C, dan Fitzpatrick17k-C.
    \item Evaluasi kinerja mencakup metrik performa klinis (Balanced Accuracy, Macro-F1, MCC), reliabilitas prediksi (Expected Calibration Error / ECE), dan efisiensi fisik perangkat keras (jumlah parameter, GFLOPs, latensi inferensi, dan konsumsi VRAM).
\end{enumerate}

\section{Manfaat Penelitian} \label{I.Manfaat}
Manfaat yang diharapkan dari hasil penelitian ini meliputi: \par

\begin{enumerate}[noitemsep]
    \item \textbf{Manfaat Akademis dan Teoritis}: memberikan bukti empiris mengenai efisiensi linier arsitektur State Space Models (Mamba) pada citra medis dermatologi, serta menyajikan metodologi evaluasi lintas-dataset yang bersih dari kebocoran data (data leakage).
    \item \textbf{Manfaat Praktis dan Klinis}: menyediakan rujukan arsitektur model diagnostik AI yang ringan, presisi, dan efisien sehingga siap diimplementasikan pada perangkat Edge AI di fasilitas kesehatan dengan keterbatasan infrastruktur.
\end{enumerate}

\section{Sistematika Penulisan} \label{I.Sistematika}
Sistematika penulisan berisi pembahasan apa yang akan ditulis disetiap Bab. Sistematika pada umumnya berupa paragraf yang setiap paragraf mencerminkan bahasan setiap Bab. \par

\subsection*{Bab I}
Bab ini berisikan penjelasan latar belakang dari topik penelitian yang berlangsung, rumusan masalah dari masalah yang dihadapi pada penjelasan di latar belakang, tujuan dari penelitian, batasan dari penelitian, manfaat dari hasil penelitian, dan sistematika penulisan tugas akhir. \par

\subsection*{Bab II}
Bab ini membahas mengenai tinjauan pustaka dari penelitan terdahulu dan dasar teori yang berkaitan dengan penelitian ini.

\subsection*{Bab III}
Bab ini berisikan penjelasan alur kerja sistem, alat dan data yang digunakan, metode yang digunakan, dan rancangan pengujian.

\subsection*{Bab IV}
Bab ini membahas hasil implementasi dan pengujian dari penelitian yang dilakukan, serta analisis dan evaluasi yang dapat dipetik dari hasil.

\subsection*{Bab V}
Bab ini membahas kesimpulan dari hasil penelitian dan juga saran untuk penelitian selanjutnya.